\scanpage{166}
\deleted{\unsure{$\pi$. Je}} $P_i$ se déduit de $P'_i$ par une
\uline{section} de $P'_i/\pi$ passant par $\overline{\xi'}$. $S$ \unsure{étant|etant}
connexe, il ne peut y en avoir qu'une
seule telle section.

c) Si $\pi\mathrel{\unsure{-|\to}}\pi'\rightrightarrows\pi''$ est exact,
alors \deleted{\unsure{$H(\,$}} $\unsure{\Phi|\varphi}(\pi)\mathrel{-}\Phi(\pi')\rightrightarrows\Phi(\pi'')$ aussi.

Il \deleted{\unsure{reste|veux}} à prouver que si $P'$ de groupe $\pi'$,
pt marqué $\xi'$, est tel que \deleted{\unsure{les}} l'on
ait un isom $P'\times_{\pi'}(\pi'',\varphi_1)\overset{u}{\xrightarrow{\sim}}P'\mathbin{\unsure{\times|*}}_{\pi'}(\pi'',\varphi_2)$,
alors on peut restreindre le groupe structural.

Or soit $Q$ le sous-schéma de $\unsure{P|P'}$
des coïncidences de $ui_1$ et $i_2$.

% diagram in the left margin, above the triangle
\diagram{\begin{tikzcd}[ampersand replacement=\&,column sep=small]
\pi \arrow[r] \& \pi' \arrow[r,shift left=.45ex] \arrow[r,shift right=.45ex] \& \pi''
\end{tikzcd}}
\diagram{\begin{tikzcd}[ampersand replacement=\&,column sep=small,row sep=small]
P' \arrow[r,"i_1"] \arrow[dr,"i_2"'] \& P''_1 \arrow[d,"u"]\\
\& P''_2
\end{tikzcd}}

Il est ouvert et fermé, contient
$\xi'$, est stable par $\pi$, donc
correspond à un sous-schéma de
$P'/\pi$ à la fois ouvert et fermé $S'$
[descente \inserted{essentiellement} \unsure{fidèlt|fidèlement} plat et \unsure{quasi-compacte|quasi-compact} \unsure{---|...}]
De plus, on voit que les \deleted{\unsure{degrés}} fibres géom
de $S'\mathrel{-}S$ sont partout un pt ; donc
c'est un isom $S'\to S$ \unsure{..}, on a gagné\unsure{...}
% archive sheet number: 165
